\section{旧接口，新名号：晚西汉到新朝的连续与断裂}
\label{sec:han-late-western-continuity-breaks}

王莽受禅不应把西汉和东汉之间切成一段没有日常生活的空白。晚西汉的尚书、人事、
外戚、户籍、钱币与边地供给仍在运作；改变的是哪些人能够解释这些接口、哪些
名称需要重新写进文书、市场和地方关系。改官名、改币制与土地政策的雄心越大，
越依赖县吏、商人、乡里与边郡反复辨认。命令被宣布并不等于同样速度地被接受，
更不等于已经改变占有、债务和保护关系。

灾荒与流动因而不能只被写成新朝失败的背景。离开土地的家庭、需要粮食的城邑、
想保住交易的商人和寻找武装保护的地方人，都会把中枢的分类改写成自己的选择。
绿林、赤眉、更始和刘秀周围的行动者并非同一社会集团；他们拥有不同的道路、
资源和名义。雒阳复汉能够逐步取得优势，不是旧汉国号自行恢复效力，而是它在
河北、关东、关中、陇右和巴蜀重新寻找能输粮、传令和安排人户的节点。

因此，度田和限制掠卖奴婢应与受禅和战争连读。前者试图使战后土地与人户重新
可见，后者试图让被战争改变的人身关系获得申诉；两者都会触动地方已经形成的
利益。新朝的断裂没有取消旧接口，反而暴露它们若失去信任与供给，名称和仪式
无法独自维持帝国。
